\chapter{Marco teórico}

\section{Tecnologías de medición de nivel}
% Categorías existentes y diferentes tecnologías
% Sensores por capacitancia

\section{Capacitancia}
% Modelo de placas paralelas
% Dependencia del área, constante dieléctrica y distancia de separación
% Modelo cilíndrico
% Modelo elíptico (Eccéntrico)
% Constante dieléctrica de los materiales (profundizar en N2O)
% Dependencia de la constante dieléctrica respecto a la densidad del N2O y, por lo tanto, de la temperatura

\section{Osciladores}
% Criterio de Barkhausen
% Osciladores de relajación (LC y RC)

\section{Fuerzas de corte o cizallamiento}
% Definición
% Ecuaciones
% Relevancia para este proyecto

\section{Conceptos de instrumentación}
% --- ENFOQUE GENERAL ---
% "Sensors and Signal Conditioning" – Ramón Pallás-Areny y John G. Webster
% "Measurement, Instrumentation, and Sensors Handbook" – editado por John G. Webster
% "Principles of Measurement Systems" – John P. Bentley

% --- METROLOGÍA ---
% "Metrología" – Carlos Sánchez Pérez / o también "Metrología y Metrotecnia" – varios autores latinoamericanos
% Guía GUM (Guide to the Expression of Uncertainty in Measurement) – BIPM - No es un libro comercial, es un documento normativo internacional, pero es LA referencia obligada si vas a hacer análisis de incertidumbre. Se puede descargar gratis del sitio del BIPM.
% NIST Technical Notes – el NIST (EE.UU.) tiene guías gratuitas sobre caracterización de sensores.